% Teil II: der durchgerechnete Weg.

\subsection*{Gegeben und gesucht}
$10!=3628800$ endet mit zwei Nullen; gesucht: mit wie vielen Nullen endet
$100!$?

\subsection*{Ansatz}
Endnullen = Primfaktoren $5$ in $100!$ (Fünfen sind der limitierende
Faktor, da Zweien reichlich vorhanden).

\subsection*{Rechnung}
\begin{align}
  &\left\lfloor\frac{100}{5}\right\rfloor = 20 &&\text{(Vielfache von 5)}\\
  &\left\lfloor\frac{100}{25}\right\rfloor = 4 &&\text{(Vielfache von 25 geben je eine Fünf extra)}\\
  &\left\lfloor\frac{100}{125}\right\rfloor = 0 .
\end{align}
Summe: $20+4=24$.

\begin{ergebnis}
$100!$ endet mit \textbf{24} Nullen.
\end{ergebnis}

\begin{kontrolle}
Schätzung: $5,10,\dots,100$ sind 20 Vielfache; davon $25,50,75,100$ je eine
zusätzliche Fünf ($4$), Summe $24$. Die $2$er-Potenz in $100!$ ist deutlich
größer ($\ge 24$), also limitiert allein die Fünf.
\end{kontrolle}

\begin{merksatz}
Endnullen von $n!$ = Anzahl der $5$er inklusive derer aus $25,125,\dots$;
Zwei Faktoren sind nie der Engpass.
\end{merksatz}
